\section{Relative excess improvement at every scale}
\label{sec:excess-propagation}

The preceding section identifies the limiting normal heights as affine
functions. We now turn that information into approximating planes.
Two details matter: the heights are integrated against changing measures,
and their normal coordinates belong to changing orthogonal frames.
Polarization first gives strong convergence to fixed affine
representatives. An explicit graph construction then gives planes with
vanishing normalized excess. A contradiction argument makes this
conclusion uniform in the measure, center, and scale.

Throughout this section, fix integers $1\le n<d$ and constants
$c,G>0$ and $D_0\ge0$. We consider Radon measures $\sigma$ satisfying
\begin{align}
 \sigma(B(x,s))&\le Gs^n
       &&(x\in\R^d,\ s>0),\label{eq:prop-upper}\\
 \sigma(B(x,s))&\ge cs^n
       &&(x\in\supp\sigma,\ 0<s\le\diam(\supp\sigma)),
       \label{eq:prop-lower}\\
 \|R_{\sigma,\eta}f\|_{L^2(\sigma)}
   &\le D_0\|f\|_{L^2(\sigma)}
       &&(\eta>0,\ f\in L^2(\sigma)).\label{eq:prop-riesz}
\end{align}
The last condition includes existence of the indicated $L^2$ output.
All truncations and pairings are those defined in Section~\ref{sec:background}.
The lower bound in \eqref{eq:prop-lower} is required only at admissible
radii; the upper bound in \eqref{eq:prop-upper} holds at every radius.

\subsection{Polarization under changing measures}

The first lemma records exactly how the moment convergence from
Lemma~\ref{lem:moving-compactness} becomes strong convergence. It uses an
actual Lipschitz function as its target.

\begin{lemma}[Polarization on a bounded region]
\label{lem:affine-polarization}
Let $\mu_j$ and $\nu$ be finite measures on $\R^d$, all concentrated on
one bounded set, and suppose that $\mu_j$ converges weakly to $\nu$.
Let $w_j\in L^2(\mu_j)$ and $f\in L^2(\nu)$. Suppose that
\begin{equation}\label{eq:polarization-norm}
 \int w_j^2\,d\mu_j\longrightarrow\int f^2\,d\nu
\end{equation}
and, for every compactly supported Lipschitz function $\psi$,
\begin{equation}\label{eq:polarization-moment}
 \int w_j\psi\,d\mu_j\longrightarrow\int f\psi\,d\nu.
\end{equation}
If $v:\R^d\to\R$ is globally Lipschitz and $f=v$ almost everywhere
for $\nu$, then
\[
 \int|w_j-v|^2\,d\mu_j\longrightarrow0.
\]
\end{lemma}

\begin{proof}
Choose a compactly supported Lipschitz cutoff equal to one on a ball
containing the common bounded set. Its product with $v$ is a compactly
supported Lipschitz function $\widetilde v$ agreeing with $v$ on that
set. Equation~\eqref{eq:polarization-moment}, with
$\psi=\widetilde v$, gives
\[
 \int w_jv\,d\mu_j\longrightarrow
 \int fv\,d\nu=\int f^2\,d\nu.
\]
Weak convergence tested against the bounded continuous function
$\widetilde v^2$ gives
\[
 \int v^2\,d\mu_j\longrightarrow
 \int v^2\,d\nu=\int f^2\,d\nu.
\]
The restrictions of $v$ belong to the respective $L^2$ spaces because
they are bounded on the common set. Thus all terms in the identity
\[
 \int|w_j-v|^2\,d\mu_j
 =\int w_j^2\,d\mu_j-2\int w_jv\,d\mu_j
       +\int v^2\,d\mu_j
\]
are integrable. Their limits prove the assertion.
\end{proof}

Let $L:\R^n\to\R^d$ be the isometry onto the limiting plane.
Use the regions $\Omega_m$ of \eqref{eq:height-projection-regions},
with coordinate map $L^*$. Each fixed ball lies in some $\Omega_m$.
Lemma~\ref{lem:moving-compactness} supplies weak convergence of the
finite restrictions and convergence of their signed and squared moments.

An intrinsic affine function $t\mapsto b+Mt$ on $L(\R^n)$ has the
globally Lipschitz representative
\begin{equation}\label{eq:ambient-affine-representative}
 v(x)=b+ML^*x.
\end{equation}
Thus Lemma~\ref{lem:affine-polarization} applies after affine rigidity.
No representative of an arbitrary limiting $L^2$ class is evaluated
against a different measure.

\subsection{Normal coordinates and tilted planes}

Write $q_0=d-n$. Let $a_j\in\R^d$, and let
$L_j:\R^n\to\R^d$ and $N_j:\R^{q_0}\to\R^d$ be linear isometries
forming a complete orthogonal splitting:
\begin{equation}\label{eq:complete-normal-splitting}
 L_j^*N_j=0,\qquad L_jL_j^*+N_jN_j^*=I_{\R^d}.
\end{equation}
For positive numbers $\delta_j$, define the normalized normal height
and the tangent projection by
\begin{equation}\label{eq:normalized-vector-height}
 w_j(x)=\delta_j^{-1}N_j^*(x-a_j),\qquad
 P_jx=a_j+L_jL_j^*(x-a_j).
\end{equation}

\begin{lemma}[Affine heights produce affine planes]
\label{lem:normal-affine-planes}
Let $\mu_j$ be measures, let the frames satisfy
\eqref{eq:complete-normal-splitting}, and suppose that
$\delta_j>0$ tends to zero. Assume
\[
 \sup_j\int\|w_j\|^2\,d\mu_j<\infty
\]
and that, for one fixed affine map $v(x)=b+Bx$ from $\R^d$ to
$\R^{q_0}$, each function $\|w_j-v\|^2$ is integrable and
\[
 \int\|w_j-v\|^2\,d\mu_j\longrightarrow0.
\]
Then there are affine $n$-planes $S'_j$ such that
\begin{equation}\label{eq:tilted-plane-distance-limit}
 \int\left(\frac{\dist(x,S'_j)}{\delta_j}\right)^2
       d\mu_j(x)\longrightarrow0.
\end{equation}
Every integral in this conclusion is finite. No convergence of the
normal frames $N_j$ is required.
\end{lemma}

\begin{proof}
The orthogonal decomposition gives
\begin{equation}\label{eq:physical-normal-error}
 x-a_j=L_jL_j^*(x-a_j)+\delta_jN_jw_j(x),
 \qquad |x-P_jx|=\delta_j\|w_j(x)\|.
\end{equation}
Since $v$ has Lipschitz constant at most $\|B\|$,
\begin{equation}\label{eq:affine-projection-replacement}
 \int\|v(x)-v(P_jx)\|^2\,d\mu_j(x)
 \le\|B\|^2\delta_j^2\int\|w_j\|^2\,d\mu_j
 \longrightarrow0.
\end{equation}
This estimate also proves integrability of the left side.
The inequality $\|u+z\|^2\le2\|u\|^2+2\|z\|^2$ now yields
\begin{equation}\label{eq:projected-affine-strong}
 \int\|w_j(x)-v(P_jx)\|^2\,d\mu_j(x)\longrightarrow0.
\end{equation}

Set $c_j=b+Ba_j$, $A_j'=BL_j$, and define
\begin{equation}\label{eq:genuine-tilted-plane}
 S'_j=\bigl\{a_j+L_jt+\delta_jN_j(c_j+A_j't):t\in\R^n\bigr\}.
\end{equation}
Its direction is the range of $L_j+\delta_jN_jA_j'$. By
\eqref{eq:complete-normal-splitting},
\[
 L_j^*(L_j+\delta_jN_jA_j')=I_{\R^n}.
\]
The parametrizing linear map is therefore injective. Its range has
dimension $n$, and \eqref{eq:genuine-tilted-plane} is a nonempty affine
$n$-plane. This conclusion needs no bound on its slope.

For a given $x$, choose $t=L_j^*(x-a_j)$ in
\eqref{eq:genuine-tilted-plane}. Since $v(P_jx)=c_j+A_j't$,
\eqref{eq:physical-normal-error} gives
\[
 \frac{\dist(x,S'_j)}{\delta_j}
 \le\|w_j(x)-v(P_jx)\|.
\]
The squared left side is measurable and dominated by an integrable
function. Integrating and using \eqref{eq:projected-affine-strong}
proves \eqref{eq:tilted-plane-distance-limit}.
\end{proof}

Lemma~\ref{lem:moving-compactness} treats the finite family of normal
coordinates on one common subsequence. Proposition~\ref{prop:height-equation}
and Lemma~\ref{lem:fractional-rigidity} identify those same limiting
functions without further extraction.

\subsection{The sequential improvement}

\begin{proposition}[Vanishing normalized excess]
\label{prop:normalized-excess-limit}
Fix $0<\gamma<1/2$. Let $\mu_j$ satisfy
\eqref{eq:prop-upper}--\eqref{eq:prop-riesz} with the same constants,
and suppose that $0\in\supp\mu_j$. Let $\delta_j>0$ and
$K_j\in\mathbb N$ satisfy
\begin{equation}\label{eq:sequential-horizon-schedule}
 \delta_j\to0,\qquad K_j\to\infty,\qquad
 \delta_j2^{\gamma K_j}\to0,\qquad
 \delta_j^{-1}2^{-K_j}\to0.
\end{equation}
Assume that $2^{K_j}\le\diam(\supp\mu_j)$ and
\begin{equation}\label{eq:sequential-horizon-excess}
 e_{\mu_j}(0,2^\ell)\le\delta_j2^{\gamma\ell}
       \qquad(0\le\ell\le K_j).
\end{equation}
Put $A_j=8\,2^{K_j}$. Suppose, for every compactly supported
1-Lipschitz scalar function $\varphi$ with
$\supp\varphi\subset B(0,A_j)$ and $\int\varphi\,d\mu_j=0$, that
\begin{equation}\label{eq:sequential-cubic-oscillation}
 |\mathcal R_{\mu_j}(\varphi)|\le\delta_j^3.
\end{equation}
There is one subsequence along which, for every fixed $r>0$,
\begin{equation}\label{eq:sequential-excess-conclusion}
 \frac{e_{\mu_j}(0,r)}{\delta_j}\longrightarrow0,
 \qquad b_{\mu_j}(0,r)\longrightarrow0.
\end{equation}
\end{proposition}

\begin{proof}
Admissibility of the growing horizons implies
$\diam(\supp\mu_j)\to\infty$.
Lemma~\ref{lem:plane-comparison} and
Proposition~\ref{prop:flat-limit}, applied to
\eqref{eq:sequential-horizon-excess}, give, after discarding finitely
many indices and passing to a subsequence, a full constant-density plane
limit
\[
 \mu_j\rightharpoonup\nu=q\H^n|_{L(\R^n)},\qquad q>0,
\]
where $L$ is a linear isometry. They also give approximate
fitting planes $S_j=a_j+L_j(\R^n)$, with $a_j\to0$ and $L_j\to L$,
and fixed-plane moment bounds
\begin{equation}\label{eq:prop-fixed-plane-moments}
 \int_{B(0,2^\ell)}\dist(x,S_j)^2\,d\mu_j(x)
 \le C_*\delta_j^2(2^\ell)^{n+2+2\gamma}
       \qquad(0\le\ell\le K_j).
\end{equation}
The constant $C_*$ is independent of $j$ and $\ell$.
The fitting planes are approximate minimizers; no existence of an exact
minimizer is used.

Choose complementary normal isometries $N_j$ so that
\eqref{eq:complete-normal-splitting} holds. For
$i=1,\dots,q_0$, write
\[
 u_{j,i}(x)=\langle N_je_i,x-a_j\rangle,
 \qquad w_{j,i}(x)=\delta_j^{-1}u_{j,i}(x),
\]
where $(e_i)$ is the standard orthonormal basis of $\R^{q_0}$.
The identity
$\sum_i u_{j,i}(x)^2=\dist(x,S_j)^2$ supplies the normal-coordinate
moment bounds required by Lemma~\ref{lem:height-energy}.
That lemma and Lemma~\ref{lem:moving-compactness}, with a common
subsequence as described above, give measurable limiting heights $f_i$
on the limiting plane. On every region $\Omega_m$, they satisfy
\begin{align}
 \int_{\Omega_m}w_{j,i}^2\,d\mu_j
   &\longrightarrow\int_{\Omega_m}f_i^2\,d\nu,
       \label{eq:normal-coordinate-norm-limit}\\
 \int_{\Omega_m}w_{j,i}\psi\,d\mu_j
   &\longrightarrow\int_{\Omega_m}f_i\psi\,d\nu
       \label{eq:normal-coordinate-signed-limit}
\end{align}
for every compactly supported Lipschitz $\psi$.
All these statements concern the same heights and the same subsequence.

Proposition~\ref{prop:height-equation} gives the mean-zero fractional
equation and the weighted tail integrability of each $f_i$.
Lemma~\ref{lem:fractional-rigidity} therefore gives a scalar $b_i$ and
a linear functional $M_i$ on $\R^n$ such that
\[
 f_i(Lt)=b_i+M_it\quad\text{for almost every }t\in\R^n.
\]
Use the ambient affine representative
$v_i(x)=b_i+M_iL^*x$. Applying
Lemma~\ref{lem:affine-polarization} to the finite restrictions on
$\Omega_m$ yields
\[
 \int_{\Omega_m}|w_{j,i}-v_i|^2\,d\mu_j\longrightarrow0.
\]
Assemble $v=(v_1,\dots,v_{q_0})=b+Bx$.
The Euclidean squared norm is the sum of the squared coordinates, so
\begin{equation}\label{eq:vector-affine-strong-limit}
 \int_{\Omega_m}\|w_j-v\|^2\,d\mu_j\longrightarrow0,
 \qquad \sup_j\int_{\Omega_m}\|w_j\|^2\,d\mu_j<\infty.
\end{equation}
Integrability of each coordinate square justifies the finite summation.

Fix $r>0$ and choose $m$ with $B(0,r)\subset\Omega_m$.
The integrands in \eqref{eq:vector-affine-strong-limit} are
nonnegative. Restricting them to $B(0,r)$ preserves the uniform bound
and the zero limit. In particular, this step requires no assertion about
mass on the sphere $\partial B(0,r)$.
Lemma~\ref{lem:normal-affine-planes}, applied to the restrictions on
this ball, supplies affine $n$-planes $S'_j$ with
\[
 \int_{B(0,r)}\left(\frac{\dist(x,S'_j)}{\delta_j}\right)^2
       d\mu_j(x)\longrightarrow0.
\]
By the definition of the excess,
\[
 0\le\left(\frac{e_{\mu_j}(0,r)}{\delta_j}\right)^2
 \le r^{-n-2}\int_{B(0,r)}
       \left(\frac{\dist(x,S'_j)}{\delta_j}\right)^2d\mu_j(x).
\]
This proves the first limit in
\eqref{eq:sequential-excess-conclusion} for every fixed $r$, without
another extraction.

Proposition~\ref{prop:flat-limit} also gives bilateral convergence of
the supports to $L(\R^n)$ on every bounded ball. This convergence
persists on the subsequence used for the height limits. For every fixed
$r>0$, both directed deviations in the definition of $b_{\mu_j}(0,r)$
therefore tend to zero when evaluated at the plane $L(\R^n)$.
Their sum divided by $r$ bounds $b_{\mu_j}(0,r)$ from above,
which proves the second limit in \eqref{eq:sequential-excess-conclusion}.
\end{proof}

\subsection{Uniform propagation at every admissible scale}

The next theorem is the input used in the cell-packing argument. The
normalization of its tests agrees with that of $W_J(Q)$ after replacing
the cell scale by $r$.

\begin{theorem}[Uniform excess propagation]
\label{thm:excess-propagation}
Fix $1\le n<d$, $c,G>0$, $D_0\ge0$, $0<\gamma<1/2$, and
$\varepsilon>0$. For every integer $J_0\ge0$, there is an integer
$J\ge J_0$, depending only on these fixed data, with the following
property. Set
\begin{equation}\label{eq:uniform-propagation-constants}
 \delta=2^{-J},\qquad K=2J+3,\qquad A=8\,2^K.
\end{equation}
Let $\mu$ satisfy \eqref{eq:prop-upper}--\eqref{eq:prop-riesz},
let $a\in\supp\mu$, and let $r>0$. Assume
\begin{align}
 2^Kr&\le\diam(\supp\mu),\label{eq:prop-admissible-horizon}\\
 e_\mu(a,2^\ell r)&\le\delta2^{\gamma\ell}
       &&(0\le\ell\le K),\label{eq:prop-horizon-assumption}
\end{align}
and assume
\begin{equation}\label{eq:prop-oscillation-assumption}
 |\mathcal R_\mu(\varphi)|\le\delta^3r^n
\end{equation}
for every compactly supported scalar function $\varphi$ satisfying
\[
 \Lip(\varphi)\le r^{-1},\qquad
 \supp\varphi\subset B(a,Ar),\qquad
 \int\varphi\,d\mu=0.
\]
Then
\begin{equation}\label{eq:uniform-propagation-conclusion}
 e_\mu(a,r/2)\le\delta,\qquad
 b_\mu(a,r/2)<\varepsilon.
\end{equation}
Equivalently, the oscillation hypothesis can be stated with
$\Lip(\varphi)\le1$ and threshold $\delta^3r^{n+1}$.
The constants $J,\delta,K,A$ are chosen before the measure, center,
and radius. The test radius $Ar$ need not be admissible.
\end{theorem}

\begin{proof}
We first work at center zero and scale one. If the assertion failed,
then for each $j\ge0$ we could choose a measure $\mu_j$ satisfying
the hypotheses at the index $j+J_0$, for which at least one conclusion
in \eqref{eq:uniform-propagation-conclusion} fails. Thus, with
\[
 \delta_j=2^{-(j+J_0)},\qquad
 K_j=2(j+J_0)+3,\qquad A_j=8\,2^{K_j},
\]
we would have either
$e_{\mu_j}(0,1/2)>\delta_j$ or
$b_{\mu_j}(0,1/2)\ge\varepsilon$ at every index.
The constants in the growth and operator hypotheses remain
$c,G,D_0$ for this entire sequence. Moreover,
\begin{align*}
 \delta_j2^{\gamma K_j}
   &=2^{3\gamma}\,2^{-(1-2\gamma)(j+J_0)}\longrightarrow0,\\
 \delta_j^{-1}2^{-K_j}
   &=2^{-(j+J_0)-3}\longrightarrow0.
\end{align*}
Thus every hypothesis of
Proposition~\ref{prop:normalized-excess-limit} holds.
On its common subsequence, the excess divided by $\delta_j$ is
eventually less than one and the bilateral beta number is eventually
less than $\varepsilon$. This contradicts the selected failure of
their conjunction. We have obtained one $J\ge J_0$ that works for
every source measure at center zero and scale one.

Fix this $J$. For arbitrary $\mu,a,r$ in the statement, define the
normalized measure
\[
 \sigma=r^{-n}(T_{a,r})_\#\mu,
 \qquad T_{a,r}(x)=\frac{x-a}{r}.
\]
The upper and admissible lower growth constants are unchanged, the
origin belongs to its support, and
$2^K\le\diam(\supp\sigma)$. The conjugation of hard truncations
under $T_{a,r}$ preserves the bound $D_0$; the truncation parameter
$\eta$ for $\sigma$ corresponds exactly to $r\eta$ for $\mu$.
Transporting all affine $n$-planes by this affine bijection gives
\begin{equation}\label{eq:prop-flatness-covariance}
 e_\sigma(0,s)=e_\mu(a,rs),\qquad
 b_\sigma(0,s)=b_\mu(a,rs).
\end{equation}
For the first equality, the squared distance contributes $r^{-2}$,
the measure contributes $r^{-n}$, and the normalizing radius contributes
$s^{-n-2}$; the result is exactly the expression normalized by
$(rs)^{n+2}$. The second equality follows directly by scaling both
directed distances and the radius by $r$.

If $\psi$ is a compactly supported 1-Lipschitz test in $B(0,A)$ with
zero $\sigma$-mean, then
$\varphi(x)=\psi(T_{a,r}x)$ has Lipschitz constant at most $r^{-1}$,
is supported in $B(a,Ar)$, and has zero $\mu$-mean.
The homogeneity of the kernel, applied to the renormalized pairing of
Lemma~\ref{lem:pairings}, gives
\[
 \mathcal R_\sigma(\psi)=r^{-n}\mathcal R_\mu(\varphi).
\]
Hence \eqref{eq:prop-oscillation-assumption} supplies the normalized
bound $|\mathcal R_\sigma(\psi)|\le\delta^3$.
The scale-one assertion applies to $\sigma$, and
\eqref{eq:prop-flatness-covariance} gives
\eqref{eq:uniform-propagation-conclusion} for $\mu$.
Finally, multiplying or dividing a test by $r$ proves the equivalence
of the two oscillation normalizations.
\end{proof}

We henceforth fix $\gamma=1/4$. The smallness parameter in
Theorem~\ref{thm:excess-propagation} can be made smaller than any
prescribed positive constant by increasing $J_0$ before choosing $J$.
The conclusion is simultaneous in excess and bilateral flatness;
there is no need to reconcile separately selected thresholds.
